% ══════════════════════════════════════════════════════════════════════════════
%  EXPERIMENT 4 — BEAKS AND FOOD                         (Lesson 8, T3 Week 3)
% ══════════════════════════════════════════════════════════════════════════════
\experimentdiv{Beaks and Food}{Modelling an adaptation \textbullet{} Fair trials \textbullet{} Correlation and cause}{8}\label{sec:4}
\lessoncodes{VC2S8U04 \textbullet{} VC2S8I01, I02, I04, I05}

\vocabhead
\begin{ctab}
\begin{tabular}{|SW{4.2cm}|W{11.6cm}|}
  \hline\hd & \thead{Word} & \thead{Definition (copy from the slide)} \\ \hline
  & \scaf{\textbf{Adaptation}} & \scaf{A feature or behaviour that helps an organism survive and reproduce in its habitat.} \\ \hline
  & \textbf{Structural adaptation} & \\ \hline
  & \textbf{Niche} & \\ \hline
  & \textbf{Correlation} & \\ \hline
  & \textbf{Cause} & \\ \hline
\end{tabular}
\end{ctab}

\begin{lipanel}
\inquiryli{Run fair repeated trials, graph the results and decide whether a pattern shows cause.}
\understandli{Explain how the shape of a beak is a structural adaptation to a bird's food.}
\sctitle
\begin{itemize}[leftmargin=1.4em, itemsep=0pt, topsep=1pt]
  \item I can keep the trials fair: same time, same food, same rules.
  \item I can draw a column graph of the averages.
  \item I can link each ``beak'' to a real Australian bird and its food.
\end{itemize}
\end{lipanel}

\begin{skillsfocus}
\focus{QP1}{Write a hypothesis for one food.}
\focus{PC2}{Keep every trial fair.}
\focus{PA1}{Record counts in a THUD table.}
\focus{PA2}{Draw a column graph.}
\focus{PA6}{Decide whether the pattern shows cause.}
\focus{PA8}{Say what the model leaves out.}
\end{skillsfocus}

% ─── TECHNIQUE ────────────────────────────────────────────────────────────────
\exppart{Technique spotlight}{Blue}

\begin{technique}{Correlation or cause?}{PA6}
Two things are \textbf{correlated} when they change together: as one goes up, the other goes up (or down). That alone does not show that one \textbf{causes} the other. You can claim cause only when you changed one thing yourself and kept everything else the same, as in a fair test.

\smallskip
\emph{Ice-cream sales and shark sightings rise together.} Correlated, but ice-cream does not cause sharks: hot weather causes both.
\end{technique}

\begin{workedex}{Correlation or cause?}
\textbf{Claim:} ``The spoon picked up more rice than the tweezers, so a spoon-shaped beak causes better rice-eating.''

\smallskip
In this experiment the beak was the only thing changed (same rice, same 30 seconds, same student), so the difference was caused by the beak shape \emph{in the model}. Whether it applies to real birds needs evidence from real birds.
\end{workedex}

\begin{wedo}{Correlation or cause?}
Students who eat breakfast get higher test marks. Does breakfast cause higher marks? Give one other thing that could explain the pattern.
\answerbox{3.5cm}
\end{wedo}

% ─── PLAN AND SAFETY / METHOD ────────────────────────────────────────────────
\exppart{Plan, safety and method}{Teal}

\textbf{Aim:} To investigate the effect of changing the shape of a ``beak'' on the amount of each food it can collect in 30 seconds.

\begin{center}
\begin{tikzpicture}[scale=0.62, font=\scriptsize]
  % tweezers
  \begin{scope}[shift={(0,0)}]
    \draw[eq] (0,0) -- (1.6,2.6) (0.35,0) -- (1.7,2.5); \draw[eq] (1.6,2.6) -- (1.7,2.5);
    \node at (0.85,-0.45) {tweezers};
  \end{scope}
  % peg
  \begin{scope}[shift={(3.0,0)}]
    \draw[eq] (0,0) rectangle (0.32,2.6); \draw[eq] (0.42,0) rectangle (0.74,2.6); \draw[eq] (-0.1,1.4) rectangle (0.84,1.8);
    \node at (0.37,-0.45) {peg};
  \end{scope}
  % spoon
  \begin{scope}[shift={(5.6,0)}]
    \draw[eq] (0.5,0.55) ellipse (0.5 and 0.55); \draw[eq] (0.5,1.1) -- (0.5,2.6);
    \node at (0.5,-0.45) {spoon};
  \end{scope}
  % straw
  \begin{scope}[shift={(8.0,0)}]
    \draw[eq] (0,0) rectangle (0.3,2.6);
    \node at (0.15,-0.45) {straw};
  \end{scope}
  % foods
  \begin{scope}[shift={(10.4,0)}]
    \draw[eq] (0,0.2) .. controls (0,-0.1) and (1.6,-0.1) .. (1.6,0.2); \draw[eq] (0,0.2) -- (1.6,0.2);
    \foreach \x in {0.3,0.6,0.9,1.2}{\fill[Ink] (\x,0.35) ellipse (0.09 and 0.04);}
    \node at (0.8,-0.45) {rice};
    \draw[eq] (2.2,0.2) .. controls (2.2,-0.1) and (3.8,-0.1) .. (3.8,0.2); \draw[eq] (2.2,0.2) -- (3.8,0.2);
    \foreach \x in {2.5,2.9,3.4}{\draw[Ink, fill=white] (\x,0.42) circle (0.17);}
    \node at (3.0,-0.45) {marbles};
    \draw[eq] (4.4,0) rectangle (5.6,1.0); \fill[black!15] (4.42,0.02) rectangle (5.58,0.7);
    \node at (5.0,-0.45) {``nectar''};
  \end{scope}
\end{tikzpicture}
\expfigcap{Four ``beaks'' and three foods.}
\end{center}

\task Choose one food. Write a hypothesis about which beak will collect the most of it. \skr{QP1}
\answerbox{2.5cm}

\needspace{12\baselineskip}
\task Complete the variables table. \skr{PC2}
\begin{ctab}
\begin{tabular}{|SW{4.4cm}|W{10.2cm}|}
  \hline\hd & \thead{Variable} & \thead{In this experiment} \\ \hline
  & \scaf{Independent} & \scaf{The beak (tweezers, peg, spoon or straw)} \\ \hline
  & Dependent & \\ \hline
  & Controlled & \\ \hline
  & Controlled & \\ \hline
\end{tabular}
\end{ctab}

\subsubsection*{Equipment}
Tweezers, clothes peg, plastic spoon, straw \textbullet{} rice grains, marbles, a cup of coloured water (``nectar'') \textbullet{} 3 shallow trays \textbullet{} a cup as the ``stomach'' \textbullet{} 10\,mL measuring cylinder \textbullet{} stopwatch

\begin{infobox}{Safety}
Marbles on the floor are a slip hazard: pick up any that fall straight away. The straw is used like a dropper, never sucked: nothing goes in your mouth in the lab.
\end{infobox}

\subsubsection*{Method}
\begin{enumerate}[itemsep=2pt]
  \item Spread 50 rice grains in one tray and 20 marbles in another. Put 50\,mL of ``nectar'' in a cup.
  \item The ``bird'' holds the beak in one hand and the stomach cup in the other. The timer says go.
  \item For 30 seconds, move as much food as possible into the stomach cup, using only the beak. For nectar, dip the straw in, cover its top with a finger, lift it out and let the nectar run into the stomach cup.
  \item Count the rice grains or marbles, or measure the nectar in mL. Reset the tray.
  \item Do three trials for each beak with each food, swapping the ``bird'' each time.
\end{enumerate}

% ─── RESULTS ─────────────────────────────────────────────────────────────────
\exppart{Results}{Amber}

\subsubsection*{Table 4.1: Average food collected in 30 seconds (three trials) \skr{PA1}}
\begin{ctab}
\begin{tabular}{|SW{3.2cm}|C{3.4cm}|C{3.4cm}|C{3.8cm}|}
  \hline\hd & \thead{Beak} & \thead{Rice (grains)} & \thead{Marbles (marbles)} & \thead{Nectar (mL)} \\ \hline
  & \scaf{Example: chopsticks} & \scaf{11.3} & \scaf{4.0} & \scaf{0.0} \\ \hline
  & Tweezers & & & \\ \hline
  & Peg & & & \\ \hline
  & Spoon & & & \\ \hline
  & Straw & & & \\ \hline
\end{tabular}
\end{ctab}
{\small\color{DarkGrey}Record each trial in your workbook or on the back page, then write the averages here to one decimal place.}

\needspace{18\baselineskip}
\task Draw a column graph of the average rice collected by each beak. Give it a title and label both axes with units. \skr{PA2}
\graphgrid{8cm}

% ─── ANALYSIS ────────────────────────────────────────────────────────────────
\exppart{Analysis and evaluation}{Coral}

\task Which beak was best for each food? Use numbers from Table 4.1. \skr{PA5}
\answerbox{4cm}

\needspace{12\baselineskip}
\task Match each beak to a real bird and its food. \skr{PA8}
\begin{ctab}
\begin{tabular}{|SW{3.0cm}|W{5.6cm}|W{6.0cm}|}
  \hline\hd & \thead{Beak} & \thead{Real bird} & \thead{Its food} \\ \hline
  & \scaf{Straw} & \scaf{New Holland honeyeater} & \scaf{Nectar from flowers such as grevillea} \\ \hline
  & Tweezers & & \\ \hline
  & Peg & & \\ \hline
  & Spoon & & \\ \hline
\end{tabular}
\end{ctab}

\needspace{12\baselineskip}
\task Did the beak shape \emph{cause} the differences in your model? Explain using how the trials were run. \skr{PA6}
\answerbox{4.5cm}

\needspace{10\baselineskip}
\task Give two ways this model is different from a real bird feeding. \skr{PA8}
\answerbox{4cm}

\needspace{10\baselineskip}
\task Write a conclusion for the food you chose: the trend, numbers that back it, and whether your hypothesis was supported. \skr{EV1}
\answerbox{4.5cm}

\begin{yr8box}
\begin{keepext}
\exttask On the Galápagos Islands, a drought killed most plants with small soft seeds. The next year, finches with deeper beaks were more common. Explain this using the idea of adaptation, and say why the finches did not ``grow'' deeper beaks on purpose.
\answerbox{5cm}
\end{keepext}
\end{yr8box}

% ─── SKILL CHECK ─────────────────────────────────────────────────────────────
\exppart{Skill check}{NavyBlue}

\begin{skillcheck}
\scrow{QP1}{I can write a hypothesis that names both variables.}
\scrow{PC2}{I kept every trial fair.}
\scrow{PA1}{My table has a title, headings, units and matching decimals.}
\scrow{PA2}{I drew a column graph with a title and labelled axes.}
\scrow{PA6}{I can tell correlation from cause.}
\scrow{PA8}{I can say what a model shows well and what it leaves out.}
\end{skillcheck}

\task Which skill will you focus on in Experiment 5? Explain what you will do differently.
\answerbox{5cm}
